% Copyright 2026 MarcosHCK
%
\documentclass{article}[12pt,a4paper]
\usepackage{amsfonts}
\usepackage{amsmath}
\usepackage{amsthm}
\usepackage{bookmark}
\usepackage{hyperref}
\usepackage{cleveref} % must load after hyperref
\usepackage{fontawesome5}
\usepackage{graphicx}
\usepackage{inputenc}[utf8]
\usepackage{indentfirst}
\usepackage{minted}
\usepackage{polyglossia}[spanish,english]
\usepackage{biblatex}[style=apa] % must load after polyglossia
\usepackage{csquotes} % must load after biblatex
\usepackage{setspace}
\usepackage{svg}
\usepackage{xcolor}[dvipsnames]

\setmainlanguage{spanish}
\setotherlanguage{english}

\setstretch{1.5}

\newcommand{\bashps}{\textasciitilde\#\,}

\definecolor{Aquamarine}{rgb}{0.498, 1, 0.83}
\definecolor{DarkAquamarine}{rgb}{0.18, 0.55, 0.34}
\definecolor{YellowOrange}{rgb}{1, 0.7, 0}
\newcommand\myshade{85}
\colorlet{mycitecolor}{YellowOrange}
\colorlet{mylinkcolor}{violet}
\colorlet{myurlcolor}{DarkAquamarine}

\hypersetup {
    colorlinks = true,
    citecolor  = mycitecolor!\myshade!black,
    linkcolor  = mylinkcolor!\myshade!black,
    urlcolor   = myurlcolor!\myshade!black,
  }

\begin{document}

\begin{titlepage}

  \centering
  {\huge\bfseries Reporte (semanas 1 y 2)\par}
	{\Large Autor: \itshape Marcos Antonio Pérez Lorenzo\par}
\end{titlepage}

\pagebreak
\section{Arquitecturas generales}

\subsection{C++}
\label{sec:genarc:cplusplus}

Todos los programas escritos en \textbf{C++} siguen las siguientes características

\begin{itemize}
  \item Lenguaje: \textbf{C++}.
  \item Dependencias: librería estándar.
  \item Flujo: obtención de datos, cálculos, salida del programa.
\end{itemize}

\subsection{Python}
\label{sec:genarc:python}

Todos los programas escritos en \textbf{Python} siguen las siguientes características

\begin{itemize}
  \item Lenguaje: \textbf{Python}.
  \item Dependencias: librería estándar.
  \item Flujo: obtención de datos, cálculos, salida del programa.
\end{itemize}

\section{Instrucciones de uso}

\subsection{C++}
\label{src:genericinstrc:cplusplus}

Todos los programas escritos en \textbf{C++} se pueden ejecutar mediante un script de \textbf{Python} intermediario para su fácil uso, por lo tanto se ejecutan de la misma forma usado para programas escritos en \textbf{Python} (descritos en \ref{src:genericinstrc:python})

\subsection{Python}
\label{src:genericinstrc:python}

Todos los programas escritos en \textbf{Python} se pueden ejecutar desde una terminal con acceso al interprete. Por ejemplo, para ejecutar el programa "clase567": en una terminal cuyo directorio de trabajo es el directorio raíz del repositorio, ejecutar "\mintinline{bash}{python3 clase567}".

\section{Programas}

Todos los programas se encuentran alojados en este \href{https://github.com/MarcosHCK/master-programming/}{repositorio \faGithubSquare}

\subsection{Programa 1}

\subsubsection{Objetivos}

Varios ejemplos de programas hechos en Python, finalmente componiendo el código para calcular el volumen de un prisma, obteniendo las longitudes de las lados del usuario a través de la función \mintinline{python3}{input (...)}.

\subsubsection{Tecnicismos}

Arquitectura basada en \ref{sec:genarc:python}.

\subsubsection{Código relevante}

El núcleo del código incluye la obtención y tratamiento de los datos, en dos líneas:

\begin{minted}{python3}
  values = ( float (input (f'Lado {n} (m): ')) for n in [ 'A', 'B', 'C' ] )
  volume = reduce (mul, values)
\end{minted}

\subsubsection{Ejecución}

Ver sección \ref{src:genericinstrc:python}.

\subsubsection{Resultados esperados}

\bashps\mintinline{bash}{python3 clase1}

\begin{itemize}
  \item[] Lado A (m): 1
  \item[] Lado B (m): 1
  \item[] Lado C (m): 1
  \item[] El volumen es 1.0m\textasciicircum3
\end{itemize}

\subsubsection{Liga de \textbf{Github}}

\href{https://github.com/MarcosHCK/master-programming/tree/master/clase1}{Aquí \faGithubSquare}

\subsection{Programa 2}

\subsubsection{Objetivos}

Crear un programa para simular un semi-sumador de 1 (un) bit

\subsubsection{Tecnicismos}

Arquitectura basada en \ref{sec:genarc:python}.

\subsubsection{Código relevante}

El núcleo del código incluye el tratamiento de los datos:

\begin{minted}{python3}
  print (f'{a} + {b} = {a ^ b} (carry {a and b})')
\end{minted}

\subsubsection{Ejecución}

Ver sección \ref{src:genericinstrc:python}.

\subsubsection{Resultados esperados}

\bashps\mintinline{bash}{python3 clase2}

\begin{itemize}
  \item[] Primer bit: 1
  \item[] Segundo bit: 0
  \item[] 1 + 0 = 1 (carry 0)
\end{itemize}

\subsubsection{Liga de \textbf{Github}}

\href{https://github.com/MarcosHCK/master-programming/tree/master/clase2}{Aquí \faGithubSquare}

\subsection{Programa 3}

\subsubsection{Objetivos}

Crear un programa para simular las puertas lógicas estudiadas en clase

\subsubsection{Tecnicismos}

Arquitectura basada en \ref{sec:genarc:python}.

\subsubsection{Código relevante}

El núcleo del código incluye el tratamiento de los datos:

\begin{minted}{python3}
  operations = [
    (lambda a, b: (a and b, a ^ b), 'suma', '(c, r)'),
    (lambda a: int (not a), 'not', 'r'),
    (lambda a, b: a and b, 'and', 'r'),
    (lambda a, b: a or b, 'or', 'r'),
    (lambda a, b: a ^ b, 'xor', 'r'),
    (lambda a, b: int (not (a and b)), 'nand', 'r'),
    (lambda a, b: int (not (a or b)), 'nor', 'r'),
    (lambda a, b: int (not (a ^ b)), 'xnor', 'r'), ]
\end{minted}

\subsubsection{Ejecución}

Ver sección \ref{src:genericinstrc:python}.

\subsubsection{Resultados esperados}

\bashps\mintinline{bash}{python3 clase3}

\begin{itemize}
  \item[] suma de (a, b) as (c, r)
  \item[] --------------------------
  \item[] suma de (0, 0) es (0, 0)
  \item[] ...
          not de (a, b) as r
  \item[] --------------------------
  \item[] not de (0,) es 1
  \item[] ...
          and de (a, b) as r
  \item[] --------------------------
  \item[] and de (0, 0) es 0
  \item[] ...
\end{itemize}

\subsubsection{Liga de \textbf{Github}}

\href{https://github.com/MarcosHCK/master-programming/tree/master/clase3}{Aquí \faGithubSquare}

\subsection{Programa 4}

\subsubsection{Objetivos}

Crear un programa que calcule el resultado de un sistema de ecuaciones lineales de primer grado con dos variables, utilizando cálculo matricial, específicamente la regla de Cramer.

\subsubsection{Tecnicismos}

Arquitectura basada en \ref{sec:genarc:python}. Utiliza la regla de Cramer para calcular los valores de ambas incógnitas.

\subsubsection{Código relevante}

El núcleo del código incluye el captura de datos (de un archivo externo):

\begin{minted}{python3}
  def read_row (line: str, n: int = 0):

    pieces = filter (lambda p: len (p) > 0, line.split (' '))
    pieces = tuple ( int (p) for p in pieces )

    if n > 0 and len (pieces) != n:
      raise Exception ('invalid matrix row length')

    return pieces

  with (Path (__file__.removesuffix ('/__main__.py')) / 'example.txt').open ('rt') as stream:

    a, b, c = read_row (stream.readline (), 3)
    d, e, f = read_row (stream.readline (), 3)
\end{minted}

Y el tratamiento de datos:

\begin{minted}{python3}
  if 0 == (di := e * a - d * b):
    raise Exception ('can not solve a singular matrix')

  y = (f * a - d * c) / di
  x = c / a - (b / a * y)
\end{minted}

\subsubsection{Ejecución}

Ver sección \ref{src:genericinstrc:python}.

\subsubsection{Resultados esperados}

\bashps\mintinline{bash}{python3 clase4}

\begin{itemize}
  \item[] Primer bit: 1
  \item[] Segundo bit: 0
  \item[] 1 + 0 = 1 (carry 0)
\end{itemize}

\subsubsection{Liga de \textbf{Github}}

\href{https://github.com/MarcosHCK/master-programming/tree/master/clase4}{Aquí \faGithubSquare}

\subsection{Programa 5}

\subsubsection{Objetivos}

Es un programa que condensada, varias clases en uno, cuyos objetivos son:

\begin{enumerate}
  \item Crear un programa que calcule el resultado de un sistema de ecuaciones lineales de primer grado con tres variables.
  \item Crear un programa que calcule el resultado de un sistema de ecuaciones lineales de primer grado con cuatro variables.
  \item Crear un programa que calcule el resultado de un sistema de ecuaciones lineales de primer grado con \(n\) variables.
\end{enumerate}

El programa logrado calcula la solución para sistema de ecuaciones lineales de primer grado con \(n\) variables, de la cual los dos primeros objetivos son un caso específico.

\subsubsection{Tecnicismos}

Arquitectura basada en \ref{sec:genarc:cplusplus}. Utiliza una forma de descomposición \textbf{LU}, llamada descomposición \textbf{LU} con pivote, donde se calculan dos matrices \(L\) y \(U\) triangulares inferior y superior, respectivamente, que satisfacen le ecuación \(PA=LU\), donde \(P\) es una matriz de permutación, y \(A\) es la matriz a resolver. Finalmente, el vector solución se calcula de forma \(Ly=Pb\) y \(Uy=x\), donde \(b\) es el vector de términos independientes, \(y\) es un vector intermedio, y \(x\) es el vector resultado; ambas ecuaciones son trivialmente solubles, dado que ambas matrices \(L\) y \(U\) son triangulares.

\subsubsection{Código relevante}

El núcleo del código incluye la captura de datos:

\begin{minted}{c++}
  static inline std::pair<std::vector<double>, std::vector<double>> load_max (std::istream& stream)
  {

    std::vector<double> row, rows;
    std::vector<double> b;
    decltype (row)::size_type last = 0;

    for (std::string line (1, 0); std::getline (stream, line);)
      {

        if (line.starts_with ("#") || line.size () == 0)
          continue;

        auto [ line_, b_i ] = split_b (line);

        b.push_back (b_i);

        std::sregex_token_iterator end;
        std::sregex_token_iterator iter (line_.begin (), line_.end (), space_rgx, -1);

        for (; iter != end; ++iter)
          {

            auto str = iter->str ();

            ltrim (line_);
            rtrim (line_);

            if (str.length () > 0)
              row.push_back (std::stod (str));
          }

        if (last == 0)
          rows.reserve (last = row.size ());

        else if (last != row.size ())
          throw wrap_stacktrace<std::runtime_error> ("incompatible amount of columns in a row");

        rows.insert (rows.end (), row.begin (), row.end ());
        row.clear ();
      }

    if (last == 0)
      throw wrap_stacktrace<std::runtime_error> ("unexpected empty matrix");

    if (last != rows.size () / last)
      throw wrap_stacktrace<std::runtime_error> ("unexpected rectangular matrix");

  return { rows, b };
  }
\end{minted}

Y el tratamiento de datos:

\begin{minted}{c++}
  static inline std::vector<double> solve (std::vector<double> A, std::vector<double>& b)
  {

    std::vector<double> p (b.size ());
    std::vector<double> y (b.size ());
    std::vector<double> x (b.size ());

    for (decltype (b.size ()) n = b.size (), i = 0; i < n; i++)
      p [i] = i;

    for (decltype (b.size ()) n = b.size (), k = 0; k < n; ++k)
      {

        auto maxRow = (decltype (b.size ())) k;
        auto maxVal = std::abs (A [k + n * p [k]]);

        for (decltype (n) i = k + 1; i < n; ++i)

          if (auto next = std::abs (A [k + n * p [i]]); maxVal < next)
            {
              maxVal = next;
              maxRow = i;
            }

        if (maxVal < 1e-15)
          {
            auto m = "Matrix is singular (maxVal " + std::to_string (maxVal) + ")";
            throw wrap_stacktrace<std::runtime_error> (m);
          }

        std::swap (p [k], p [maxRow]);

        for (decltype (n) i = k + 1; i < n; ++i)
          {

            auto factor = A [k + n * p [i]] / A [k + n * p[k]];

            A [k + n * p [i]] = factor;
            
            for (decltype (n) j = k + 1; j < n; ++j)
              A [j + n * p [i]] -= factor * A [j + n * p [k]];
          }
      }

    for (decltype (b.size ()) n = b.size (), i = 0; i < n; ++i)
      {

        double sum = 0.;

        for (size_t j = 0; j < i; j++)
          sum += A [j + n * p[i]] * y [j];

        y [i] = b [p[i]] - sum;
      }

    for (decltype (b.size ()) n = b.size (), i = n - 1; i >= 0; --i)
      {

        double sum = 0.;

        for (decltype (n) j = i + 1; j < n; j++)
          sum += A [j + n * p[i]] * x [j];

        x [i] = (y [i] - sum) / A [i + n * p[i]];
      }
  return x;
  }
\end{minted}

\subsubsection{Ejecución}

Ver sección \ref{src:genericinstrc:cplusplus}.

\subsubsection{Resultados esperados}

\bashps\mintinline{bash}{python3 clase567}

\begin{itemize}
  \item[] x = [ 0.0382182, 0.000944757, -0.0709671 ], took 4us
\end{itemize}

\subsubsection{Liga de \textbf{Github}}

\href{https://github.com/MarcosHCK/master-programming/tree/master/clase567}{Aquí \faGithubSquare}

\end{document}